\section{Donoso on Liberalism, Socialism, and Catholicism}

Although this book by Cortes was not one of Evola's favourites because of its overt theological emphasis, it is still very valuable reading for several audiences. First of all, for counter-revolutionaries, it makes clear the ideas of liberalism and socialism and brings to awareness that this opposition to Tradition is nothing recent, but has been ongoing for a few hundred years. For anti-Catholics, it will make clear the Catholicism at one time was the sole serious and viable force opposed to the modern world. They need to ask themselves if they are really willing to reject the world that defined Europe, a world created by foundational documents of Augustine, Aquinas, Dante, or Bossuet, not to mention art, architecture, chivalry and so much else. For neo-Catholics in thrall to the ``spirit of Vatican II'', they need to decide if they are in continuity with the Church described by Donoso. Since the legitimacy and authority of the Catholic Church derives only from its adherence to Tradition, neo-Catholics have to demonstrate they are still Catholic.

Donoso's Essays on Liberalism, Socialism, and Catholicism with a new introduction by Cologero Salvo is now available in the Library\footnote{\url{http://www.gornahoor.net/library/CortesEssays.pdf}}.

\begin{quotex}
My principles are only those that, before the French Revolution, every well-born person considered healthy, sane and normal.
\flright{\textsc{Julius Evola}}
\end{quotex}


In this short list of modern authors---that is, those writing after 1789---Evola includes Juan Donoso Cortes among those holding to those sane and normal principles. He writes:

\begin{quotex}
In the same spirit ... of the great Catholic philosophers of authority, Joseph de Maistre and Donoso Cortes, I deny everything that, directly or indirectly, derives from the French revolution, which, in my opinion, has Bolshevism as its ultimate outcome, in contrast to the World of Tradition.
\end{quotex}


To the modern mind, such a view is summarily dismissed as ``Fascism'', the common accusation hurled at any idea opposed to the ``manifest destiny'' of the modern world. Yet as Evola points out, this spirit of Tradition predates the Fascist movements of the 20th century because:


\begin{quotex}
It brings back a tradition higher and prior to Fascism, insofar as it belongs to the heritage of a hierarchical, aristocratic, and traditional conception of the state, conceptions having a universal character that were maintained in Europe right up until the French Revolution.
\end{quotex}


In contrast to the totalitarianisms of the modern age, this traditional view was based on an organic society, with power distributed under the principle of subsidiarity, and bonds of loyalty freely given.

In these essays, Donoso reaches down to the fundamental assumptions of the modern world as expressed in Liberalism and Socialism, in contrast to the spirit of Tradition that Catholicism used to represent. Since few men ever bother to question and articulate clearly their assumptions, a common dialog among the three world views becomes literally impossible. Donoso writes:

\begin{quotex}
The supreme interest of that school is in preventing the arrival of the day of radical negations or of sovereign affirmations; and that it may not arrive, it confounds by means of discussion all notions, and propagates scepticism, knowing as it does, that a people which perpetually hears in the mouth of its sophists the pro and the contra of everything, ends by not knowing which side to take, and by asking itself whether truth and error, injustice and justice, stupidity and honesty, are things opposed among themselves, or are only the same things regarded from different points of view.
\end{quotex}


How much more true this is in our day, with instant and incessant televised discussions and the multitude of Internet blogs. The intent is seldom to discern truth, justice, and intelligence, but instead the goal is to promote a particular point of view.

\paragraph{Liberalism and Socialism}

This is how well-born men used to think:

\begin{quotationx}
The Liberal school regards it as certain, that there is no other evil but what is in the political institutions which we have inherited from time to time, and that the supreme good consists in levelling those institutions in the dust. The greater part of Socialists look upon it as settled, that there is no other evil but what is in society, and that the grand remedy is in the complete destruction of social institutions. All agree that the evil comes to us from times past; the Liberals affirm that the good can be realised at the present time, and the Socialists that the golden age cannot commence till the future.

The supreme good consisting, according to one and the other, in a supreme disarrangement, which, according to the Liberal school, must be realised in the political, and according to the Socialists, in the social, regions, the one and the other agree in the substantial and intrinsic goodness of man, who is to be the intelligent and free agent of both disarrangements. This conclusion has been explicitly enunciated by the Socialistic schools, and is implicitly involved in the theory which the Liberal schools hold. That conclusion follows from their theory, in such a way, that if the conclusion be denied, the theory itself comes to the ground. In fact, the theory, according to which the evil is in man and proceeds from man, is contradictory of the other, according to which, the evil is in the social or political institutions, and proceeds from the political and social institutions. Supposing the former, what logically follows is, to extirpate the evil in man, with which its extirpation in society and in government must necessarily be secured. Supposing the latter, what logically follows is, to extirpate the evil directly in society or in government, in which are its centre and its origin.

From which we see that the Catholic, and the rationalistic, theories are not only incompatible, but even contradictory. By the Catholic, disturbance, whether political or social, is condemned as mad and useless. The rationalistic theories condemn all moral reform of man as useless and mad. And the one and the other are consistent in their condemnations; for if the evil be not in government or in society, why will you disturb society or government? And on the contrary, if the evil is not in individuals, nor proceeds from individuals, why will you attempt an interior reform of man?
\flright{\textsc{Donoso Cortes}}
\end{quotationx}


In a nutshell, Donoso Cortes points out the fundamental worldview of the three schools. Since no one ever reaches down to their fundamental assumptions, dialogue is literally impossible. There can only be incessant and fruitless debate. As Donoso puts it:

\begin{quotex}
The supreme interest of that school is in preventing the arrival of the day of radical negations or of sovereign affirmations; and that it may not arrive, it confounds by means of discussion all notions, and propagates scepticism, knowing as it does, that a people which perpetually hears in the mouth of its sophists the pro and the contra of everything, ends by not knowing which side to take, and by asking itself whether truth and error, injustice and justice, stupidity and honesty, are things opposed among themselves, or are only the same things regarded from different points of view.
\end{quotex}


So we are forced to take one or the other ``side'' of the debate, when what is really necessary is to understand one's fundamental assumptions. And many so-called conservatives hold a reactionary position on political or social issues, more or less out of old habit, while their fundamental worldview differs little from liberals or socialists. To summarize, see Table \ref{tab:PositionsSources}.

\begin{table}[h]\small\centering
\begin{tabular}{cc}\toprule

\textbf{Position} & \textbf{Source of Evil} \\\midrule

Liberalism & Political institutions\\ 

Socialism & Social institutions \\

Tradition & Man himself \\\bottomrule
\end{tabular}
\caption{Positions and Sources of Evil}
\label{tab:PositionsSources}
\end{table}

From this chart we can see where each school is logically compelled to find a solution:

\begin{itemize}


\item Liberals look to change political institutions

\item Socialists look to change social mores

\item Tradition looks to the moral regeneration of man
\end{itemize}


As Donoso points out, the liberals are ultimately incoherent. Their political institutions are lifeless and without soul. The liberal regards them as self-existent apart from the peoples, their values, history and consciousness. The conservative senses this but is unable to articulate it properly, primarily because he is a believer himself in his political institutions.

The socialist believes that the solution lies in the overthrow of traditional societal norms; only then will justice and happiness be achieved. Here we see the constant push against and challenge to traditional morality and social norms.

So we hear the cry from the conservative that the battle must be fought in the social arena. However, he himself absorbs the socialist view on issues --- if not now, then eventually --- due to the constant flow of propaganda. Hence, he gradually concedes the ground that he had initially won.

The war is really about worldview but no one notices. Metaphysics must be opposed to the rationalism of the liberals and socialists; the moral regeneration of man must be given primary place before political and social change. These voices are few, and where they do exist, they are unconvincing and their philosophy is inconsistent.


\flrightit{Posted on 2010-03-12 by Cologero}

\begin{center}* * *\end{center}

\begin{footnotesize}
\begin{sffamily}
\texttt{Cologero on 2010-03-23 at 06:50 said:}

... the Liberal school has done nothing but establish premises which end in Socialistic consequences, and the Socialistic schools, nothing but draw the consequences contained in the Liberal premises. Those two schools differ not in ideas, but in daring: when the question thus stands between them, it is clear the victory belongs by right to the more daring, and the more daring without any doubt is that which, without stopping midway, accepts their consequences with the principles. If this be so, there is no doubt ... that Socialism has the best of the battle, and that hers is the palm of victory.
\flright{\textsc{Donoso Cortes}}


\hfill

\texttt{Attila on 2011-12-22 at 23:37 said:}

I tend toward the view that ``Man'' IS the ``problem'' ---even though I have no connection with Catholicism. All of the secular -isms have left people without any reference points -- except those furnished by the media complex.

\hfill

\texttt{Charlotte Cowell on 2011-12-23 at 06:20 said:}

There is an incredible all-round dumbing down of people and it only seems to be getting worse. Although on the one hand I see more and more people unfolding into their spiritual awakening -- in accordance with the divine plan -- it is quite astonishing how clueless human beings are becoming on the other. the distinction between the haves and have nots is not one of money but of wisdom-knowledge. Education to put it simply. Of course this is a ploy -- possibly unconscious -- on the part of governments and other institutions to keep people stupid, to feed them bread at the circuses so they remain bloated and content without the desire to seek for more. When I worked in a large office I was routinely shocked by the lack of true education I saw among people -- young people especially -- who are actually considered to be top graduates and recruits.The lack of understanding was not simply spiritual but historical and political, it made me wonder what on earth they had been doing to have escaped university with first class degrees but zero wisdom-knowledge. It made me very sad actually and I wondered what hope there is, but of course hope comes from the spiritual hierarchies who will individually awaken (the memories of) people with the revelatory fire of divine love, the only means by which true wisdom-knowledge may be attained. But you are right to say the media has a lot to answer for, so many of those corporations are in a prime position to offer wonderful educations that will help open people's minds, but they instead take the quick road to greed by filling their heads with reality TV nonsense and tabloid news. One is therefore beholden to remind one's self of who owns these corporations... ..we know the answer, the old new world order has poor taste in entertainment, the new new world order should take over ASAP :-)

\end{sffamily}
\end{footnotesize}

